%----------------------------------------------------------------------------
\chapter{Bevezető}%\addcontentsline{toc}{chapter}{Bevezető}
%----------------------------------------------------------------------------

Sok nyelvtanuló szembesül azzal a problémával, hogy bár hosszú időn keresztül tanul egy idegen nyelvet, mégsem érzi magát magabiztosnak, amikor valódi beszédhelyzetbe kerül. Gyakran előfordul, hogy az elsajátított tudása ellenére nehézséget jelent egy egyszerű párbeszéd lebonyolítása is, vagy a gondolatok megfogalmazása. Sokszor hallhatjuk azt is, hogy egy nyelvet akkor lehet igazán megtanulni, ha azt aktívan használjuk\cite{swain1985communicative,long1996role}, azonban erre nem minden tanulónak van lehetősége. Különösen igaz ez azokra, akiknek a környezetükben nincs olyan személy, akivel rendszeresen gyakorolhatnák a nyelvet.

A hagyományos nyelvtanulási módszerek és alkalmazások jellemzően nyelvtani szabályokra, elszigetelt szókincs elsajátítására helyezik a hangsúlyt. Bár ezek fontos elemei a nyelvtanulásnak, gyakran kevés figyelem jut a beszédkészség fejlesztésére és a valós kommunikáció gyakorlására. \cite{swain1985communicative,long1996role}

Ezekre a problémákra szeretnék megoldást nyújtani az általam fejlesztett alkalmazással, amely egy mesterséges intelligenciával támogatott beszédalapú nyelvtanulási környezetet biztosít. Az alkalmazás lehetőséget nyújt arra, hogy a felhasználók mesterséges intelligenciával folytatott párbeszédek során gyakorolhatják a célnyelvet, mintha egy valódi beszélgetőpartnerrel kommunikálnának. \cite{xiao2024conversational} Az AI nemcsak részt vesz a beszélgetésben, hanem felismeri és kijavítja a nyelvi hibákat \cite{lyster2013oral,sato2012peer}, illetve eltárolja azokat, későbbi feldolgozás céljából.

A rendszerben külön figyelem irányul arra, hogy a felhasználók ne csupán javítást kapjanak, hanem lehetőséget arra is, hogy saját hibáikból tanuljanak. A gyakran előforduló hibák ismételten megjelennek a tanulási folyamat, beszélgetések és különböző gyakorlási módszerek segítségével, így elősegítve a helyes nyelvi minták rögzülését és a beszédkészség folyamatos fejlődését.

A következő fejezetekben bemutatom az alkalmazás tervezésének és megvalósításának folyamatát, az alkalmazott módszereket és technológiai megoldásokat, valamint elemzem, hogy a beszélgetésközpontú, hibaalapú megközelítés miként járul hozzá a hatékony nyelvtanuláshoz és a beszédkészség fejlesztéséhez.